\section{Immutable Parent Pointers}
\label{sec:immutable-parent-pointers}

The immutable parent pointer is the foundational structural primitive of the Recursive Spawning Protocol. It is the single mechanism that makes the delegation chain a runtime-verifiable artifact rather than a forensic reconstruction, that makes chain integrity a structural guarantee rather than a policy assertion, and that renders all modes of autonomous drift architecturally impossible regardless of operational urgency, architectural complexity, or adversarial sophistication of any machine actor in the chain. This section specifies the parent pointer formally, defines the chain traversal operator and establishes the base case for the root human anchor, specifies the immutability enforcement mechanism through the Fact Layer write protocol, describes the Purpose Layer's spawn authorization role, and provides the complete chain traversal algorithm with correctness arguments.

% ------------------------------------------------------------------
\subsection{Formal Definition: ParentPointer}
\label{sec:parent-pointer-definition}

\begin{dfn}[Parent Pointer]
\label{dfn:parent-pointer}
Let $\mathcal{N}$ be the set of all agent nodes in a \bxthree{} system implementing the Recursive Spawning Protocol, and let $P$ denote the set of all Purpose Layer actors. The \emph{parent pointer} $\alpha(n)$ is defined as a total function:

\begin{equation}
\alpha : \mathcal{N} \longrightarrow P \cup \{\bot\}
\label{eq:parent-pointer-function}
\end{equation}

where $\alpha(n) = p$ if $p \in P$ is the Purpose Layer actor that authorized the spawn event resulting in node $n$'s activation, and $\alpha(n) = \bot$ if and only if $n$ is the root human anchor established at system initialization. The pointer is established \emph{exactly once} at node provisioning via a Fact Layer atomic write operation and satisfies the following immutability property for all nodes $n$ and all actors $p$:

\begin{equation}
\bigl(\alpha(n) = p \;\ \text{provisioned}(n)\bigr) \;\ \neg\exists\; \text{modify}\bigl(\alpha(n)\bigr)
\label{eq:parent-pointer-immutability}
\end{equation}

No actor --- including the Purpose Layer after provisioning, the Bounds Engine, any administrative process, or any external principal --- may execute an operation that overwrites, redirects, or deletes $\alpha(n)$ after the provisioning write is confirmed. The Fact Layer write protocol rejects any such attempt at the protocol level, before execution reaches the ledger.
\end{dfn}

Four structural constraints govern every parent pointer:

\begin{enumerate}
\item[(P1) \textbf{Uniqueness}] For every $c \in \mathcal{N}$ spawned under RSP, there exists exactly one $p \in P$ such that $\alpha(c) = p$ holds at all times after provisioning. No node may have multiple simultaneous parent pointers, and no parent pointer may be removed or redirected after establishment.

\item[(P2) \textbf{Immutability}] After the provisioning write confirming $\alpha(c) = p$ is recorded in the Fact Layer ledger, no operation exists --- from any source, with any privilege level, under any system condition --- that can modify or delete $\alpha(c)$. The pointer is permanent for the operational lifetime of $c$. This is not a policy enforced by access control; it is a structural property of the protocol. The modification operation is not defined.

\item[(P3) \textbf{Directionality}] The parent pointer is an intrinsic property of the child node, not a reference held by the parent node. The child $c$ carries $\alpha(c)$ as an immutable attribute in its own Fact Layer state record. The parent $p$ holds no mutable reference to $c$. This means the chain is traversable from any node outward to its complete lineage regardless of the current operational state of any other node in the chain. A parent that has terminated, crashed, or been decommissioned does not affect the child's ability to reference its parent through $\alpha(c)$.

\item[(P4) \textbf{Termination at Human Anchor}] The root of any RSP chain is a human principal node $h$ for which $\alpha(h) = \bot$ holds --- the null parent marker indicating that $h$ is the ultimate delegation origin. No machine actor in the RSP model may serve as a chain root. The root human anchor is designated $\hmannchor{h}$.
\end{enumerate}

The parent pointer is not a conventional pointer variable subject to access control restrictions. It is a Fact Layer state property with no defined modification operation after provisioning. The distinction is architectural: an access-controlled pointer can be modified by any principal with sufficient privileges, and immutability is enforced by the access control policy until that policy is circumvented --- through credential theft, privilege escalation, or administrative compromise. The Fact Layer write protocol makes parent pointer immutability a structural property; there is no modification pathway to circumvent.

\bigskip
\noindent\textbf{Notation.} We write $\alpha(n) = p$ to denote that $p$ is the parent of $n$. We use $p = \parent(n)$ and $n \in \children(p)$ interchangeably. The notation $\alpha(n) = \bot$ denotes that node $n$ has no parent --- it is the root human anchor. The notation $\alpha^{(k)}(n)$ denotes the $k$-fold application of $\alpha$ (e.g., $\alpha^{(2)}(n) = \alpha(\alpha(n))$).

\begin{dfn}[Spawn Event]
\label{dfn:spawn-event}
A \emph{spawn event} is the 5-tuple

\begin{equation}
\operatorname{spawn}(c,\ p,\ \sigma_c,\ t,\ \phi)
\label{eq:spawn-event}
\end{equation}

where $c$ is the newly created child node, $p \in P$ is the authorizing Purpose Layer actor, $\sigma_c$ is the authorized execution scope, $t$ is the wall-clock timestamp, and $\phi$ is the cryptographic signature produced by the Purpose Layer over the spawn request. The Fact Layer atomically: (1) creates node $c$, (2) sets $\alpha(c) \leftarrow p$, (3) sets the node's authorized scope to $\sigma_c$, (4) seals the memory envelope, and (5) appends the event to the forensic ledger $\mathcal{L}$.
\end{dfn}

The initialization of $\alpha(c)$ occurs atomically during the spawn event. This separation of authorization (Purpose Layer) from assignment (Fact Layer) is essential: it prevents a parent from claiming to have spawned a node whose parent pointer was set to a different actor, and it prevents a child from falsifying its own ancestry after provisioning.

\begin{prop}[Immutability of ParentPointer]
\label{prop:parent-pointer-immutable}
For any child node $c$ with parent $p$, once the spawn event $\operatorname{spawn}(c, p, \sigma_c, t, \phi)$ has been recorded in the forensic ledger $\mathcal{L}$, the value $\alpha(c)$ is fixed for the operational lifetime of $c$. No subsequent operation --- including operations issued by $p$, by any Purpose Layer actor, by the Bounds Engine, or by $c$ itself --- can alter $\alpha(c)$.
\end{prop}

% ------------------------------------------------------------------
\subsection{The Chain Operator and Root Human Anchor}
\label{sec:chain-operator}

The chain operator $\Chain{\cdot}$ is the primary tool for chain traversal, verification, and accountability attribution. It constructs the complete delegation lineage from any node to the human anchor by recursively applying the parent pointer function.

\begin{dfn}[Chain Operator]
\label{def:chain-operator}
For any agent node $a \in \mathcal{N}$ in a Recursive Spawning chain, the \emph{chain} $\Chain{a}$ is defined recursively as:

\begin{equation}
\Chain{a} =
\begin{cases}
\{\, a \,\} \cup \Chain{\alpha(a)} & \text{if } \alpha(a) \neq \bot \\[10pt]
\{\, a,\ h(a) \,\}                  & \text{if } \alpha(a) = \bot \quad \text{(root human anchor)}
\end{cases}
\label{eq:chain-recursive}
\end{equation}

Equivalently, expanding the recursion to reveal the full lineage:

\begin{equation}
\Chain{a} = \bigl\{\, a,\ \alpha(a),\ \alpha^{(2)}(a),\ \alpha^{(3)}(a),\ \ldots,\ \alpha^{(k)}(a) \,\bigr\}
\label{eq:chain-expanded}
\end{equation}

where $k$ is the depth of node $a$ in the chain (the number of successive parent applications required to reach the root), $\alpha^{(j)}(a)$ denotes the $j$-fold application of $\alpha$, and $\alpha^{(k)}(a)$ is the unique root node $n_0$ satisfying $\alpha(n_0) = \bot$ and $\hmannchor{n_0}$.
\end{dfn}

The chain $\Chain{a}$ is the ordered set of all nodes on the delegation path from $a$ to and including the human anchor. It includes the terminal human anchor node, which is the only node in any RSP chain with $\alpha(n) = \bot$. The chain is always finite: by the RSP termination property, every RSP chain has depth at most $D_{\max}$, so the recursion in Equation~\ref{eq:chain-recursive} terminates at the root within at most $D_{\max} + 1$ successive applications of $\alpha$.

\bigskip
\noindent\textbf{Base case: root human anchor.} The root of every RSP chain is a Purpose Layer entity designated as a human principal, established at system initialization when a human principal $h$ authorizes the first agent node $n_0$. The Purpose Layer records $h(n_0) = h$ and $\alpha(n_0) = \bot$ in the Fact Layer authorization ledger. This is the only valid chain-root in RSP: no machine actor may serve as a chain origin, and no node may have $\alpha(n) = \bot$ unless it carries the $\hmannchor{\cdot}$ designation. The non-collapsing human anchor property ensures $h(n) = h(n_0)$ for all nodes in all chains spawned from $n_0$.

The chain operator satisfies three critical properties by construction:

\begin{enumerate}
\item[(C1) \textbf{Uniqueness}] For any node $a$, $\Chain{a}$ is a singleton set constructed by iterating $\alpha$ from $a$ to the root. There is exactly one such set; no alternative chain exists for any node in a correctly provisioned RSP system.

\item[(C2) \textbf{Determinism}] $\Chain{a}$ is identical at any runtime point to what it was at the moment of $a$'s provisioning. Because $\alpha(c)$ is immutable for all $c$, no retroactive chain modification is possible. The chain is not reconstructed from logs --- it is directly computable by iterating $\alpha$ from any node until reaching $\bot$.

\item[(C3) \textbf{Human termination}] Every $\Chain{a}$ terminates at a human principal. The RSP protocol does not permit machine actors to serve as chain roots. This property is structurally enforced by the Fact Layer pre-commit hook, which rejects any spawn event whose root is not a registered human principal.
\end{enumerate}

% ------------------------------------------------------------------
\subsection{Immutability Enforcement: Fact Layer Write Protocol}
\label{sec:fact-layer-write-protocol}

The immutability of $\alpha(n)$ is not a policy assertion enforced by privilege-gated access control. It is a structural property enforced by the Fact Layer write protocol, which defines no valid modification operation for $\alpha(n)$ after provisioning. The distinction is architectural, not a matter of degree: access control enforces immutability by restricting which principals may issue modification operations; if sufficient privilege is acquired, the barrier is circumvented. The Fact Layer write protocol makes immutability the only valid state transition.

The Fact Layer enforces immutability through four properties of its write protocol:

\begin{enumerate}
\item[(W1) \textbf{No modification operation defined.}] The Fact Layer ledger interface defines exactly two operations affecting $\alpha(n)$: an initial write during provisioning and a read for audit or traversal. No modification operation is defined in the interface specification. Any request framed as a modification is rejected before reaching the ledger state machine. The operation does not exist in the protocol's valid operation set.

\item[(W2) \textbf{Pre-commit hooks.} Before any write is committed to the ledger, the Fact Layer pre-commit hook inspects the operation type. If the operation targets the $\alpha$ field of an already-provisioned node, the write is rejected. The pre-commit hook is a structural filter, not a policy. It is part of the ledger state machine's definition, not a configurable rule.

\item[(W3) \textbf{Sealed memory envelope.} The parent pointer field $\alpha(n)$ is encrypted in the node's persistent memory envelope using a Fact Layer-only symmetric key. Even a compromised Bounds Engine cannot read $\alpha(n)$ because it never has access to the decryption key. The envelope carries an HMAC-SHA-256 signature computed over its contents using a Fact Layer-only key. Any modification to the ciphertext invalidates the signature. The Fact Layer decrypts the envelope only when reading $\alpha(n)$ to service a chain-traversal or audit request; the plaintext is never exposed to the Bounds Engine or the child node's own code. After each authorized read, the Fact Layer re-encrypts and re-signs the envelope, preventing replay attacks.

\item[(W4) \textbf{Atomic provisioning.} The parent pointer $\alpha(c)$ and its Purpose Layer warrant $\phi$ are created as a single atomic Fact Layer write. There is no temporal window in which the pointer exists without authorization, or in which authorization exists without a corresponding pointer. This eliminates any possibility of a time-of-check/time-of-use inconsistency between warrant issuance and pointer assignment.
\end{enumerate}

These four properties together constitute what we call \emph{protocol-level immutability}: the immutability of $\alpha(n)$ is a property of the protocol's operation definitions, not a property of the access control policy layered over a conventional read/write interface.

The precise failure mode eliminated by the Fact Layer write protocol is retroactive manipulation of chain topology for the purpose of accountability laundering. Without protocol-level immutability, an adversarial actor could re-parent a node to a favorable parent after the fact, obscuring the original authorization chain and making drift detection impossible. With protocol-level immutability, this manipulation is structurally impossible.

% ------------------------------------------------------------------
\subsection{Spawn Authorization: Purpose Layer Role}
\label{sec:purpose-layer-spawn-authorization}

The Purpose Layer plays two structurally necessary roles in the parent pointer lifecycle: it originates the spawn authorization policy that defines the conditions under which a new parent pointer may be created, and it issues the spawn warrant that is the prerequisite for the Fact Layer provisioning write. These roles are inseparable from the parent pointer's immutability: if the Purpose Layer could issue warrants after the fact, or if a machine actor could create a parent pointer without a Purpose Layer warrant, the chain's integrity guarantees would be void.

\medskip
\noindent\textbf{Authorization policy origination.} At system initialization, the Purpose Layer defines the spawn authorization policy $P_{\mathrm{spawn}}$ and records it in the Fact Layer authorization ledger. $P_{\mathrm{spawn}}$ specifies the mandatory conditions for any valid spawn event:

\begin{enumerate}
\item[(S1) \textbf{Scope refinement.}] The child scope must be a strict subset of the parent scope: $R(c) \subset R(p)$. No lateral expansion, no superset, no equality. This ensures that each spawn event narrows the delegation scope, preserving the intent-refinement property of the chain.

\item[(S2) \textbf{Operational necessity.}] The parent must declare, in the Fact Layer authorization ledger, the precise condition requiring child spawning. The declaration must identify why the condition cannot be resolved within $R(p)$ and why the child is the minimal appropriate response. The necessity declaration is subject to human anchor review.

\item[(S3) \textbf{Depth within bound.}] The resulting chain depth must not exceed $D_{\max}$. The depth bound is Purpose-Layer-defined and Fact-Layer-enforced at spawn time.
\end{enumerate}

These three conditions are the only authorized paths to parent pointer creation. There is no bypass, no emergency override, and no operational urgency path that permits parent pointer creation without satisfying all three conditions.

\medskip
\noindent\textbf{Spawn warrant issuance.} When a Purpose Layer actor $p \in P$ authorizes a spawn event $\operatorname{spawn}(c, p, \sigma_c, t, \phi)$, it produces the cryptographic warrant $\phi$, signing the spawn request using a Purpose-Layer-only private key. The warrant $\phi$ encodes:

\begin{enumerate}
\item The identity of the child node $c$,
\item The declared parent $p$,
\item The authorized scope $\sigma_c$,
\item The timestamp $t$, and
\item A hash of the operational necessity declaration.
\end{enumerate}

The warrant $\phi$ is theFact Layer's precondition for the provisioning write: the Fact Layer will not set $\alpha(c) = p$ unless a valid $\phi$ is present in the request. The Purpose Layer's private key is held in aPurpose-Layer-onlyHSM; the Bounds Engine and all machine actors are excluded from this key by architectural separation. This means a machine actor cannot fabricate a Purpose Layer warrant, and cannot bypass the warrant requirement by claiming emergency authority.

The warrant is not merely an audit trail. It is a structural precondition that gates the only valid mechanism for creating a new parent pointer. Without a valid $\phi$, the Fact Layer provisioning write is rejected and $\alpha(c)$ is never set. The warrant and the pointer are created atomically in a single Fact Layer transaction.

% ------------------------------------------------------------------
\subsection{Chain Traversal Algorithm}
\label{sec:chain-traversal-algorithm}

The chain traversal algorithm is the runtime procedure for constructing the delegation chain from any node to its human anchor. It is provided here in full algorithmic detail with correctness arguments.

\medskip
\noindent\textbf{Algorithm: Chain-Traverse}

\begin{algorithm}
\caption{Chain-Traverse($n$) --- Build the delegation chain from node $n$ to the human anchor}
\label{alg:chain-traverse}
\begin{algorithmic}[1]
\REQUIRE $n$ is a provisioned RSP node with $\alpha$ defined
\ENSURE $\Chain{n}$ --- the ordered set of all ancestor nodes including the human anchor
\STATE $chain \leftarrow \langle\ \rangle$ \COMMENT{Initialize empty ordered list}
\STATE $current \leftarrow n$ \COMMENT{Start at the query node}

\WHILE{$\text{true}$}
    \STATE $\text{append}(chain,\ current)$ \COMMENT{Add current node to chain}
    \IF{$\alpha(current) = \bot$}
        \STATE \textbf{break} \COMMENT{Root human anchor reached --- chain complete}
    \ENDIF
    \STATE $current \leftarrow \alpha(current)$ \COMMENT{Move to parent}
\ENDWHILE

\RETURN $chain$
\end{algorithmic}
\end{algorithm}

\bigskip
\noindent\textbf{Correctness arguments.}

\textit{Termination.} The algorithm terminates because RSP chains are depth-bounded by construction (proved in Section~\ref{sec:rs-correctness}). Each iteration of the while-loop moves $\text{current}$ to the parent node $\alpha(\text{current})$, which is strictly closer to the root along the chain. Since the chain has finite depth at most $D_{\max} + 1$, the loop executes at most $D_{\max} + 1$ iterations before $\alpha(\text{current}) = \bot$ is reached. The loop therefore terminates in finite time. This argument holds regardless of the current operational state of any node in the chain, including nodes that may have terminated or entered error states, because the traversal only reads the immutable $\alpha$ values.

\textit{Correctness of the returned set.} By Definition~\ref{def:chain-operator}, $\Chain{a}$ is the set containing $a$, $\alpha(a)$, $\alpha(\alpha(a))$, and so on, until the root. The algorithm produces exactly this sequence: it starts at $a$, appends each node, moves to its parent via $\alpha$, and stops when $\alpha(\text{current}) = \bot$ is reached. The output is identical to $\Chain{a}$ by construction.

\textit{Uniqueness of result.} The algorithm produces a unique result for any input node because $\alpha$ is a total function (defined for every provisioned node) and is injective on the chain direction (every node has exactly one parent, and the root has no parent). There is no branching, no alternative path, and no non-deterministic choice at any step.

\textit{Immutability guarantee.} The algorithm reads $\alpha$ values but never writes them. The traversal is read-only and does not interact with the Fact Layer write protocol. The algorithm can be executed at any runtime point without affecting the immutability of the chain it traverses, and without creating any modification opportunities for adversarial actors.

\bigskip
\noindent\textbf{Algorithm: Chain-Verify}

The chain traversal algorithm enables Chain-Verify, which confirms the structural validity of a delegation chain at any runtime point. This algorithm is the operational counterpart to the chain operator definition: it provides the runtime verification procedure that confirms a given node's chain satisfies all RSP structural constraints.

\begin{algorithm}
\caption{Chain-Verify($n$) --- Verify the authorization chain from node $n$ to human anchor}
\label{alg:chain-verify}
\begin{algorithmic}[1]
\REQUIRE $n$ is a provisioned RSP node
\ENSURE $\text{Boolean}$ indicating whether the chain is a structurally valid RSP chain

\STATE $chain \leftarrow \text{Chain-Traverse}(n)$
\STATE $m \leftarrow \text{len}(chain)$

\STATE \COMMENT{\textit{Step V1: Verify chain terminates at human anchor}}
\IF{$\alpha(chain[m-1]) \neq \bot$}
    \STATE \RETURN $\text{false}$ \COMMENT{Root is not $\bot$ --- malformed chain}
\ENDIF
\IF{$\neg h(chain[m-1])$}
    \STATE \RETURN $\text{false}$ \COMMENT{Root is not a designated human principal}
\ENDIF

\STATE \COMMENT{\textit{Step V2: Verify each spawn event has a Purpose Layer warrant}}
\FOR{$i \leftarrow 0$ \TO $m-2$}
    \STATE $child \leftarrow chain[i]$
    \STATE $parent \leftarrow chain[i+1]$
    \IF{$\neg \text{warrant\_exists}(child,\ parent)$}
        \STATE \RETURN $\text{false}$ \COMMENT{Missing Purpose Layer warrant for spawn}
    \ENDIF
\ENDFOR

\STATE \COMMENT{\textit{Step V3: Verify scope refinement at each link}}
\FOR{$i \leftarrow 0$ \TO $m-2$}
    \STATE $child \leftarrow chain[i]$
    \STATE $parent \leftarrow chain[i+1]$
    \IF{$\neg (R(child) \subset R(parent))$}
        \STATE \RETURN $\text{false}$ \COMMENT{Scope refinement constraint violated}
    \ENDIF
\ENDFOR

\STATE \COMMENT{\textit{Step V4: Verify depth bound compliance at each node}}
\FOR{$i \leftarrow 0$ \TO $m-1$}
    \IF{$\text{depth}(chain[i]) > D_{\max}$}
        \STATE \RETURN $\text{false}$ \COMMENT{Depth bound exceeded at node}
    \ENDIF
\ENDFOR

\RETURN $\text{true}$
\end{algorithmic}
\end{algorithm}

Chain-Verify returns \texttt{true} if and only if all four verification conditions hold simultaneously. The verification inspects structural chain properties rather than current operational state. Each condition corresponds to a fundamental RSP constraint: V1 confirms the chain terminates at a human anchor, V2 confirms every spawn was authorized by the Purpose Layer, V3 confirms scope refinement was respected at every link, and V4 confirms the depth bound was respected throughout. A chain that passes Chain-Verify is, by the RSP definitions, a fully authorized delegation chain terminating at a human principal.

The algorithm is deterministic and produces the same result regardless of which node in the chain initiates verification. There is no self-serving interpretation of authorization that a drifted or orphaned node could produce, because Chain-Verify inspects the complete chain topology and warrants in the Fact Layer ledger, not the node's own claims about its authorization. Chain verification is a property of the chain topology, not a property of the node being verified.

\bigskip
\noindent\textbf{Computational cost.} Chain-Traverse runs in $O(d)$ time where $d$ is the chain depth (at most $D_{\max} + 1$). Chain-Verify runs in $O(m)$ time for the traversal plus $O(m)$ for each of the three verification loops, yielding $O(m)$ total where $m$ is chain length (also bounded by $D_{\max} + 1$). The algorithms are linear in chain length and do not depend on the total number of nodes in the system.

\bigskip
\noindent\textbf{Constant-time escalation.} The depth bound $D_{\max}$ provides a constant-time upper bound on chain traversal cost. The worst-case number of pointer dereferences required to reach the human anchor from any node is $D_{\max} + 1$, independent of the total number of nodes in the system. A system with 10,000 spawned nodes and $D_{\max} = 5$ requires at most six pointer dereferences to reach a human --- the same as a system with 10 nodes. Mutable parent pointers would allow the chain to be extended retroactively, making the effective depth potentially unbounded and the worst-case escalation time unbounded as well. Immutability is what makes the depth bound a structural guarantee, not merely a policy.

In the Bailout Protocol context (Section~\ref{sec:bailout-protocol}), chain traversal is used to identify the correct Human Accountability Anchor when an agent at any depth encounters an unresolvable exception. The upward propagation of the exception signal follows the same parent-pointer structure as chain traversal, but uses the asynchronous trigger pathway rather than the synchronous verification pathway. Chain traversal ensures that this walk never diverges, never shortcuts, and never fails to reach a human.

% ------------------------------------------------------------------
\subsection{Summary: Why Immutability Is Structural, Not Conventional}
\label{sec:immutability-structural-summary}

The immutable parent pointer is the load-bearing constraint of the Recursive Spawning Protocol. Without it, the chain is mutable --- subject to retroactive modification by actors with sufficient privilege --- and the RSP's correctness properties (drift prevention, chain integrity, termination) collapse to policy conventions rather than structural guarantees.

The immutability is architectural --- not a stronger access control policy --- because it is enforced by the Fact Layer write protocol, which defines no valid modification operation for $\alpha(n)$ after provisioning. This is distinct from access-control-based immutability in four ways:

\begin{enumerate}
\item[(1)] \textbf{No override path.} There is no privileged actor, administrative role, or emergency pathway capable of modifying $\alpha(n)$. Access control systems have override paths; the Fact Layer write protocol does not.

\item[(2)] \textbf{Failure-state resilience.} The immutability holds under all system failure conditions, including those that would disable other enforcement mechanisms. The write protocol is a local constraint that does not depend on distributed agreement.

\item[(3)] \textbf{Source-equivalence of enforcement.} All modification requests --- from any source --- receive identical treatment: protocol-level rejection. The Fact Layer does not distinguish between principals by privilege level.

\item[(4)] \textbf{Atomic warrant-pointer creation.} The parent pointer and its Purpose Layer warrant are created as a single atomic Fact Layer write. There is no temporal window in which the pointer exists without authorization, or in which authorization exists without a corresponding pointer.
\end{enumerate}

The chain operator $\Chain{a}$ is uniquely determined by the immutable parent pointers in the chain. The chain traversal algorithm computes this operator correctly and terminates. Chain-Verify provides the runtime verification procedure for chain structural validity. Together, these mechanisms make the delegation chain an immutable, verifiable, and auditable runtime artifact --- the foundation on which all other RSP properties depend.
